\begin{definition}[Carnot group in exponential coordinates: the group data]
A \emph{Carnot group datum} consists of three integers $n,m,r$ with $r\ge1$ together with
\begin{itemize}
\item a map $*:\mathbb R^n\times\mathbb R^n\to\mathbb R^n$ (the group law),
\item a continuous linear map $\pi_1:\mathbb R^n\to\mathbb R^m$ (the projection onto the
  first layer),
\item a map $\exp_{\mathrm h}:\mathbb R^m\to\mathbb R^n$ (the horizontal exponential),
\end{itemize}
subject to the following conditions.
\begin{enumerate}
\item $(\mathbb R^n,*)$ is a group whose neutral element is the origin and whose inversion
  is $x\mapsto-x$; explicitly, $(x*y)*z=x*(y*z)$, $x*0=x$, $0*x=x$, $x*(-x)=0$ and
  $(-x)*x=0$, for all $x,y,z\in\mathbb R^n$.
\item The group law is $C^\infty$ on $\mathbb R^n\times\mathbb R^n$.
\item $\pi_1(x*y)=\pi_1(x)+\pi_1(y)$ for all $x,y\in\mathbb R^n$.
\item $\exp_{\mathrm h}$ is continuous and $\pi_1(\exp_{\mathrm h}(\zeta))=\zeta$ for every
  $\zeta\in\mathbb R^m$.
\item If $r=1$ then the group law is commutative: $x*y=y*x$ for all $x,y$.
\end{enumerate}

These are the facts about a Carnot group $\mathbb G=(\mathbb R^n,*)$ of step $r$ with $m$
generators, read in the exponential coordinates
$x=\exp\big(\sum_{j\le r}\sum_{i\le m_j}x_{i,j}X_{i,j}\big)$ of Section Section~2.4,
that the rest of the development uses at the group level. Item (1) and item (2) are the
statement recorded there that in these coordinates the group law is polynomial, hence
smooth, the neutral element is $0$ and $x^{-1}=-x$; writing the inverse as $-x$ rather than
carrying it as separate data is a choice of coordinates, not an extra hypothesis. Item (3)
is the identity $\pi_1(x*y)=\pi_1(x)+\pi_1(y)$: $\pi_1$ is a homomorphism onto the
abelianization $\mathbb G/[\mathbb G,\mathbb G]\cong\mathbb R^m$, and it is linear because
it is a coordinate projection. In item (4), $\exp_{\mathrm h}(\zeta)=\exp(\sum_i\zeta_iX_i)$
is the point whose first-layer coordinates are $\zeta$ and whose higher-layer coordinates
vanish, so $\pi_1\circ\exp_{\mathrm h}=\mathrm{id}$; in particular $\pi_1$ is surjective.
Item (5) records that a Carnot group of step $1$ has $\mathfrak g=\mathfrak g_1$ with
$[\mathfrak g_1,\mathfrak g_1]=\{0\}$, so it is the Euclidean additive group.

The formalization note (item 3 of the note after Theorem Theorem~1.3) permits the
structural facts of Section Section~2.4 to be built into the definition of the setting,
provided every Carnot group in those coordinates satisfies the definition. Each condition
above is one of those facts, so the proviso holds; in particular the Euclidean additive
group on $\mathbb R^m$, with $n=m$, $r=1$, $x*y=x+y$, $\pi_1=\mathrm{id}$ and
$\exp_{\mathrm h}=\mathrm{id}$, satisfies all five.
\end{definition}
